\documentclass[]{article}
\usepackage{amsmath, amssymb, amsthm}

\usepackage{multicol, enumitem, physics, changepage}

\renewcommand{\thesubsection}{\thesection.\alph{subsection}}

\newcommand{\mrow}[3]{\left[
	\begin{array}{ccc}
		#1 & #2 & #3
	\end{array}
	\right]}
\newcommand{\mcol}[3]{\left[
	\begin{array}{c}
		#1\\
		#2\\
		#3
	\end{array}
	\right]}

\newenvironment{augmx}[1]{%
	\left[\begin{array}{@{}*{#1}{c}|c@{}}
	}{%
	\end{array}\right]
}

\newcommand{\extaugmx}[2]{
	\left[\begin{matrix}
		#1
	\end{matrix}\right|\left.\begin{matrix}
		#2
	\end{matrix}\right]
}

\newcommand{\xlrightarrow}{\xrightarrow{\hspace{0.8cm}}}
\newcommand{\reals}{\mathbb{R}}

\newcommand{\linears}{\mathcal{L}}

\newcommand{\bmat}[1]{\begin{bmatrix}#1\end{bmatrix}}

\newcommand{\letmatrix}[2]{\mathcal{M}_{#1 \times #2}(\reals)}

\newcommand{\detmat}[1]{
	\left|\hspace{0.5em}
	\begin{matrix}
		#1
	\end{matrix}
	\hspace{0.5em}\right|
}

\begin{document}
	\begin{center}
		{\Large\textbf{MATH545: Problem Set 8}}\\
		\large Grey Cole\\
		April 27, 2026
	\end{center}
	
	\section{}
	We begin with the assertion that $\mathbb{M}_n(\reals)$ is a vector space, defined previously with matrix addition and scalar multiplication.\\
	First we will show that $T(A) = A^T$ is a linear transformation.\\
	Let $A, B \in \mathbb{M}_n(\reals)$.
	\[ T(A + B) = (A + B)^T = A^T + B^T = T(A) + T(B) \]
	which satisfies additivity.\\
	Let $A \in \mathbb{M}_n(\reals)$ and $c \in \reals$.
	\[ T(cA) = (cA)^T = c \cdot A^T = c \cdot T(A) \]
	which satisfies scalar multiplication, proving $T$ is a linear transformation.\\
	Then, to further show that $T$ is an isomorphism of vector spaces, we must show that $T$ is bijective.\\
	\textbf{One-to-One:} Let $A, B \in \mathbb{M}_n(\reals)$ where $T(A) = T(B)$. Then $A^T = B^T$. If we take the transpose of both sides, we get $A = B$. We have shown $T(A) = T(B) \implies A = B$, therefore $T$ is one-to-one/injective.\\
	\textbf{Onto:} Let $A, C \in \mathbb{M}_n(\reals)$ and choose $A$ such that $A = C^T$. Then $T(A) = A^T = (C^T)^T = C$. We have shown $\forall C \in \mathbb{M}_n(\reals), \exists A \in \mathbb{M}_n(\reals)$ such that $T(A) = C$, which makes $T$ onto/surjective.\\
	Since we have shown that $T$ is injective and surjective, it is bijective, which completes the proof that it is an isomorphism of vector spaces.
	
	\section{}
	As a reminder, we define $\beta$ as the standard basis for $\mathcal{P}_2$, so $\beta = \{1, x, x^2\}$.\\
	We can compute $T$ on each basis vector of $\beta$ to find each column of $[T]_{\beta}$.
	\begin{align*}
		T(1) &= 1 + x(0) = 1\\
		T(x) &= x + x(1) = 2x\\
		T(x^2) &= x^2 + x(2x) = 3x^2
	\end{align*}
	Now we can express each result in $\beta$:
	\[
		[T(1)]_{\beta} = \bmat{1\\0\\0} \qquad
		[T(x)]_{\beta} = \bmat{0\\2\\0}\qquad
		[T(x^2)]_{\beta} = \bmat{0\\0\\3}
	\]
	These are each the columns of the matrix:
	\[ \boxed{[T]_{\beta} = \bmat{1&0&0\\0&2&0\\0&0&3}} \]
	
	\section{}
	\subsection{}
	\begin{gather*}
		\begin{aligned}
			T(1) &= x(0) + (0) = 0\\
			\\
			T(x) &= x(1) + (0) = x\\
			\\
			T(x^2) &= x(2x) + (2) = 2x^2 + 2
		\end{aligned}
		\qquad
		\begin{aligned}
			[T(1)]_{\beta_1} = \bmat{0\\0\\0}\\
			[T(x)]_{\beta_1} = \bmat{0\\1\\0}\\
			[T(x^2)]_{\beta_1} = \bmat{2\\0\\2}
		\end{aligned}
		\qquad
		\boxed{[T]_{\beta_1} = \bmat{0&0&2\\0&1&0\\0&0&2}}
	\end{gather*}
	\subsection{}
	\begin{gather*}
		\begin{aligned}
			T(1) &= x(0) + (0) = 0\\\\
			T(x) &= x(1) + (0) = x\\\\
			T(1 + x^2) &= x(2x) + (2) = 2x^2 + 2\\&= 2(1 + x^2)
		\end{aligned}
		\qquad
		\begin{aligned}
			[T(1)]_{\beta_2} = \bmat{0\\0\\0}\\
			[T(x)]_{\beta_2} = \bmat{0\\1\\0}\\
			[T(1 + x^2)]_{\beta_2} = \bmat{0\\0\\2}
		\end{aligned}
		\qquad
		\boxed{[T]_{\beta_2} = \bmat{0&0&0\\0&1&0\\0&0&2}}
	\end{gather*}
	\subsection{}
	To find $P$, we write each vector of $\beta_1$ in terms of $\beta_2$ and use the results as the columns.
	\begin{gather*}
		\begin{aligned}
			1 = 1 \cdot 1 + 0 \cdot x + 0 \cdot (1 + x^2)\\\\
			x = 0 \cdot 1 + 1 \cdot x + 0 \cdot (1 + x^2)\\\\
			x^2 = -1 \cdot 1 + 0 \cdot x + 1 \cdot (1 + x^2)
		\end{aligned}
		\qquad
		\begin{aligned}
			[1]_{\beta_2} &= \bmat{1\\0\\0}\\
			[x]_{\beta_2} &= \bmat{0\\1\\0}\\
			[x^2]_{\beta_2} &= \bmat{-1\\0\\1}
		\end{aligned}
		\qquad
		P = [I]_{\beta_2}^{\beta_1} = \bmat{1&0&-1\\0&1&0\\0&0&1}
	\end{gather*}
	\subsection{}
	To demonstrate that $[T]_{\beta_1}$ and $[T]_{\beta_2}$ are similar, we will use the equality \[[T]_{\beta_2} = P^{-1} \cdot [T]_{\beta_1} \cdot P\]
	\begin{align*}
		\bmat{0&0&0\\0&1&0\\0&0&2} &= \bmat{1&0&1\\0&1&0\\0&0&1}\bmat{0&0&2\\0&1&0\\0&0&2}\bmat{1&0&-1\\0&1&0\\0&0&1}\\
		&= UNFINISHED FIX PLS
	\end{align*}

\end{document}